\section{Related Work}
\label{sec: related_work}
% \shortparagraph{Reinforcement Learning with Verifiable Rewards.}  Recent large-scale reasoning models such as OpenAI o1~\citep{openai2024learning}, DeepSeek‑R1~\citep{guo2025deepseek} have demonstrated that reinforcement-learning-based post‑training can substantially enhance LLM reasoning. Motivated by this reinforcement learning with verifiable rewards (RLVR) (e.g.~\citet{shao2024deepseekmath,guo2025deepseek,lambert2024tulu,yang2025qwen3,hu2025open,yu2025dapo,guan2025rstar,zeng2025simplerl} among many other works) has become a major approach for post‑training LLMs to improve reasoning. A broad literature studies how to make RLVR training effective and efficient at scale, including novel reinforcement learning algorithms and objectives~\citep{yu2025dapo,liu2025understanding,yue2025vapo,zheng2025group}, verifier architecture and reward deigns~\citep{zuo2025ttrl,zhao2025learning,agarwal2025unreasonable,prabhudesai2025maximizing}, and mechanisms that manage exploration diversity and entropy~\citep{cheng2025reasoning,chen2025pass,cui2025entropy,wang2025beyond}, or takes a critical view on the current evaluation~\citep{yue2025does,zhao2025echo,hochlehnert2025sober}. Despite steady progress, a fundamental challenge is to balance exploration and exploitation along long reasoning trajectories without brittle heuristics. We adopt a simple yet effective annealed sampling schedule that front‑loads exploration and cools later steps during rollout to encourage exploration while keeping training stable.

% \shortparagraph{Exploration Control in RLVR.} A line of works that is close to our work studies how to control exploration and sampling while doing RLVR~\citep{hou2025t1, cheng2025reasoning,tangoptimizing,chen2025pass,cui2025entropy,wang2025beyond,xiong2025minimalist,deng2025trial,xu2025not,zheng2025act,dou2025improving, li2025treepo}. In particular, \citet{cheng2025reasoning} propose per-token entropy to focus exploration at branching tokens in sampling; \citet{tangoptimizing, chen2025pass} transform per-prompt rewards to optimize pass@$k$, guiding exploration across samples; \citet{cui2025entropy} control high-covariance tokens to prevent entropy collapse and sustain exploration; \citet{wang2025beyond} update only high-entropy tokens, concentrating exploration where decisions split. While more nuanced multi-threaded exploration strategies~\citep{pan2025learning} could benefit training, adapting them to the training loop for large models often introduces significant computational overhead and the challenges of maintaining massive parallelization. Different from prior work, this paper first analyze the impacts of sampling temperature and then introduce a purely sampling-level annealed schedule that encourages exploration and then progressively stabilize answers, thereby enabling discovery of new solutions while yielding more stable training.

\shortparagraph{Exploration in RLVR.} 
Reinforcement learning with verifiable rewards (RLVR)~\citep{shao2024deepseekmath,guo2025deepseek,lambert2024tulu,yu2025dapo,yang2025qwen3,hu2025open,guan2025rstar,zeng2025simplerl} has become a major paradigm for improving LLM reasoning, with a growing literature on training algorithms~\citep{yu2025dapo,liu2025understanding,yue2025vapo,zheng2025group}, reward design~\citep{zuo2025ttrl,zhao2025learning,agarwal2025unreasonable,prabhudesai2025maximizing}, and evaluation methodology~\citep{yue2025does,zhao2025echo,hochlehnert2025sober}. A central and open challenge is balancing exploration and exploitation along long reasoning trajectories~\citep{cheng2025reasoning,cui2025entropy,wang2025beyond,deng2025trial}. Recent work addresses this through diverse mechanisms: per-token entropy-guided sampling~\citep{cheng2025reasoning}, pass@$k$-targeted training objectives~\citep{tangoptimizing, chen2025pass}, high-covariance token control to prevent entropy collapse~\citep{cui2025entropy}, selective updating of high-entropy minority tokens~\citep{wang2025beyond}, and tree-based or selective rollout strategies~\citep{li2025treepo,xu2025not,zheng2025act,dou2025improving}. However, these approaches often require tracking additional statistics or multi-threaded generation~\citep{pan2025learning}, adding computational overhead. In contrast, this paper presents the first systematic analysis of sampling temperature scheduling for RLVR and introduces a purely sampling-level annealed schedule that front-loads exploration and progressively stabilizes generation, requiring zero additional computation.


The temperature annealing design of \alg is reminiscent of \emph{Simulated Annealing} (SA)~\citep{kirkpatrick1983optimization, bertsimas1993simulated}, which also uses a cooling temperature to transition from exploration to exploitation. Prior SA-inspired approaches in ML invariably apply a uniform temperature across all positions in a generated sequence~\citep{liu2021simulated, zhang2024edt, israel2025enabling, manvi2024adaptive}. 
% This approach fails to account for the heterogeneous roles of tokens at different positions \citep{wang2025beyond}. 
Our work departs from this convention by introducing an \textbf{intra-sequence annealed temperature} schedule, where the temperature varies dynamically within the generation of a single sequence. This novel approach allows for more nuanced control over exploration and leads to significant performance gains in RLVR. We provide a more detailed review of SA in 
% and its connections to our work in 
\cref{app:simulated_annealing}.